\section{From statistical identifiability to resource closure}
\label{sec:resource-map}

\subsection{The RQIR handoff}

RQIR separates statements that are often compressed into the single word ``detectable.'' Paper I asks whether a declared finite calibration leaves an exact physical source direction on which a specified ordered response differs. Paper II asks whether the detector-facing tangent generated by that direction survives nuisance profiling and covariance whitening. Its local data-only information is positive only when the target score has a component outside the nuisance span. Paper II deliberately stops before translating that surviving information into shots, spectral densities, coherence, integration time, geometry, or control effort.

The present paper takes that handoff literally. It does not modify the source hierarchy of RQIR I or the nuisance geometry of RQIR II. Instead it asks:
\begin{quote}
\emph{Given a target direction that is statistically identifiable for at least one declared experimental design, what physical or experimental resource is required to reach a specified local information target, and how should that resource be allocated across available settings?}
\end{quote}

This is an optimal experimental-design problem, but with one RQIR-specific discipline: nuisance profiling must occur \emph{inside} the resource objective. Raw Fisher information or signal-to-noise accumulated in a channel is not credited if the same channel accumulates an equally aligned nuisance direction.

\subsection{Setting-wise resource model}

Let $k=1,\ldots,K$ label declared detector settings, bands, branches, pulse sequences, baselines, or other repeatable data-acquisition modes. For one unit of an additive resource in setting $k$, let
\begin{equation}
 \bm s_k\in\mathbb R^{m_k},
 \qquad
 J_k\in\mathbb R^{m_k\times p}
\end{equation}
denote the whitened target score and nuisance-score matrix at a fixed local reference point. Let $e_k\ge0$ be the amount of scalable resource assigned to that setting. Independent repetition or stationary integration is represented by the usual square-root score scaling,
\begin{equation}
 \bm s_k\rightarrow \sqrt{e_k}\,\bm s_k,
 \qquad
 J_k\rightarrow \sqrt{e_k}\,J_k.
\end{equation}
A fixed independent nuisance constraint is represented by a positive-semidefinite precision matrix $\Lambda\succeq0$. It is intentionally \emph{not} multiplied by the science exposure unless the physical calibration process itself is included as a separately costed resource.

Stacking all active settings and using the variational identity of RQIR II gives
\begin{equation}
\boxed{
\Ires(\bm e;\Lambda)=
\min_{\bm a}
\left[
\sum_{k=1}^K e_k
 \|\bm s_k-J_k\bm a\|^2
+\bm a^T\Lambda\bm a
\right].
}
\label{eq:resource-variational}
\end{equation}
Equation~\eqref{eq:resource-variational} is the central bridge of RQIR III. It converts a resource vector into nuisance-profiled local information without treating nuisance information as free sensitivity.

\subsection{Resource geometry}

\paragraph*{Proposition 1 (resource monotonicity, concavity, and homogeneity).}
For fixed $\{\bm s_k,J_k\}$ and fixed $\Lambda\succeq0$, $\Ires(\bm e;\Lambda)$ is componentwise nondecreasing and concave on $\mathbb R_+^K$. If $\Lambda=0$, it is additionally positively homogeneous:
\begin{equation}
 \Ires(\alpha\bm e;0)=\alpha\Ires(\bm e;0),
 \qquad \alpha\ge0.
\label{eq:homogeneity}
\end{equation}

\emph{Proof.}
For fixed $\bm a$, the bracket in Eq.~\eqref{eq:resource-variational} is affine in $\bm e$, with nonnegative coefficient
$q_k(\bm a)=\|\bm s_k-J_k\bm a\|^2$ for every component. Increasing any $e_k$ therefore increases the objective for every $\bm a$, so its minimum cannot decrease. A pointwise infimum of affine functions is concave, proving concavity. When $\Lambda=0$, the entire objective is linear in a common factor $\alpha$, which may be taken outside the minimization, proving Eq.~\eqref{eq:homogeneity}. \hfill$\square$

Several consequences are immediate. First, ``more of the same data'' can increase information only when the selected design already possesses a nonzero profiled direction. Second, diminishing marginal returns can arise from self-calibration even though every raw Fisher block scales linearly with exposure. Third, a fixed external prior and a scalable science exposure obey different resource laws; treating them as one common exposure would erase the data-only/prior-assisted distinction established in Paper II.

\paragraph*{Corollary 1 (exact shared-alignment obstruction).}
If $\Lambda=0$ and there exists a single nuisance coefficient vector $\bm a_0$ such that
\begin{equation}
 \bm s_k=J_k\bm a_0
\label{eq:shared-alignment}
\end{equation}
for every setting with $e_k>0$, then
\begin{equation}
 \Ires(\bm e;0)=0
\end{equation}
for any finite or infinite common rescaling of those allocations.

This is the resource-level form of the exposure obstruction in RQIR II. A larger budget cannot repair a design family whose active target scores remain exactly reproducible by one shared nuisance displacement. The remedy is a new response direction, a physically justified independent nuisance constraint, or a change of experimental geometry.

\subsection{Information targets are not yet discoveries}

A local precision target $\sigma_\star$ can be staged by requiring
\begin{equation}
 \Ires\ge I_\star,
 \qquad
 I_\star=\sigma_\star^{-2},
\label{eq:precision-target}
\end{equation}
when the Cram\'er--Rao approximation and its regularity conditions are appropriate. Likewise, a local asymptotic Wald significance target $Z_\star$ for a specified nonzero amplitude $\beta_\star$ corresponds to
\begin{equation}
 I_\star=\frac{Z_\star^2}{\beta_\star^2}.
\label{eq:wald-target}
\end{equation}
These equations define resource-design targets, not guaranteed finite-sample coverage or discovery claims. Boundaries, nonlinear response, non-Gaussian likelihoods, look-elsewhere effects, or disconnected nuisance solutions require the full likelihood treatment already flagged at the end of RQIR II.
